\textbf{問2} $\sqrt{2}$ が無理数であることの証明．

$\sqrt{2}$ が有理数だと仮定して矛盾を導く．
$\sqrt{2} = \frac{p}{q}$（$p, q$ は自然数）とおく．既約分数にはしない．
互いに素を使うより，素因数の個数を数えた方が仮定が少なくて済むので．
両辺を 2 乗して分母を払うと
\[
  p^2 = 2q^2
\]
となる．
左辺 $p^2$ の素因数分解に現れる 2 の個数は偶数である．
右辺 $2q^2$ の素因数分解に現れる 2 の個数は奇数である．
素因数分解はただ一通りなので，これは矛盾である．
よって $\sqrt{2}$ は無理数である． $\blacksquare$
